\documentclass[11pt]{article}

\usepackage[T1]{fontenc}
\usepackage[margin=1in]{geometry}
\usepackage{amsmath,amssymb,mathtools}
\usepackage{array}

\newcommand{\Val}{\mathsf{Val}}
\newcommand{\Some}{\mathsf{some}}
\newcommand{\None}{\mathsf{none}}
\newcommand{\Obs}{\mathsf{Obs}}
\newcommand{\Law}{\mathsf{Law}}
\newcommand{\supp}{\mathop{\mathrm{supp}}}
\newcommand{\pub}{\mathsf{pub}}
\newcommand{\raw}{\mathsf{raw}}
\newcommand{\eff}{\mathsf{eff}}

\title{An Ordered Commitment Protocol: Mathematical Boundary and Proof Plan}
\author{}
\date{}

\begin{document}
\maketitle

\paragraph{Status.}
This note specifies a proposed target semantics and the mathematical obligations
for compiling the full typed source language.  Statements explicitly labelled
\emph{checked} summarize facts already represented in the Lean development.
Statements labelled \emph{proposed} are proof targets, not claims that the
present implementation has proved a whole-program compiler theorem.  In
particular, the operational interface below contains local capabilities; it
does not contain a global execution or deviation theorem as a field.

\section{The source boundary}

\paragraph{Checked source facts.}
A source program has heterogeneous public chance, guarded private commitment,
deterministic reveal, and return.  At a commitment, the owner chooses a value
in the guard subtype.  Thus legality is established in the source view at the
commitment step, not for the first time when the value is revealed.  A reveal
copies the already committed value and has no source-level failure branch.
Public chance is sampled from a distribution computed from public state.

Consequently, in the current source semantics an invalid target opening and a
failure to open are not already two presentations of a source quit.  One may
instead want a different source abstraction in which commitment records an
unchecked binding and reveal performs validation and can resolve failure.  That
is a legitimate design alternative, but it requires an explicit change to the
source transition system and its policy information: guard timing, descendant
choices, public chance, and the legal quitting execution must then all be
defined against that deferred-validation state.  This note does not silently
identify that proposed semantics with the checked current semantics.

The existing source quitting predicate compares a legal terminal quitting
execution with a supported unilateral continuation against fixed source
opponents.  The executions must have the same public source prefix immediately
before the designated source decision.  It is deliberately a predicate on
source executions: a target timeout has no incentive meaning until a separate
correspondence proof constructs the required legal source execution.

\paragraph{Checked conditional-opening mechanism.}
The source library also represents a conditional opening as a fresh optional
commitment followed immediately by its reveal.  If the original sealed value
is $x$, the fresh value is either $\None$ or $\Some(x)$.  The original secret
remains unchanged.  Choosing $\None$ is legal in every source view.  This is a
genuine source-declared quitting operation and is fundamentally different from
rewriting the original committed value after a failed reveal.

\section{Proposed ordered target model}

Let the lowered program have sites $0,\ldots,n-1$ in canonical program-counter
order.  Site $i$ has a value type $A_i$ and is one of
\[
  \mathsf{chance}(K_i),\qquad
  \mathsf{commit}(p_i,G_i),\qquad
  \mathsf{reveal}(o_i).
\]
Here $o_i$ names either an initial private field or an earlier commitment.
Values crossing the transport boundary are tagged members of the heterogeneous
sum $\sum_i A_i$.  The public code contains types, owners, declared reads,
guards, chance kernels, and reveal origins, but not the initial private values.

A runtime state should distinguish at least
\[
 (pc,B,E,D,C,H).
\]
$B$ is immutable bound material: typed candidate meanings and any exact source
inputs authenticated at binding time.  $E$ is the effective materialized store
used by later source operations.  $D$ contains values actually disclosed to
the public.  $C$ caches chance samples, and $H$ is transport and observation
history.  Conflating $B$, $E$, and $D$ makes timeout reasoning unsound: binding
a value, using a value in the effective source context, and publishing a value
are three different events.

The program counter advances only after the current operation is realized or
after an explicitly authorized resolution for that operation.  Messages may
still be submitted, delivered, replayed, malformed, mistyped, or out of order;
such actions can affect transport-visible history, but they do not silently
realize a source step.

\subsection{Local operational capabilities}

The following capabilities are local and source-independent enough to belong
at the runtime boundary.

\begin{enumerate}
\item A typed decoder checks the site tag and produces an $A_i$, or rejects.
\item One sufficient ideal alternative is a \emph{certified recoverable
      binding} for an ordinary commitment.  It hides
      a definite value $v:A_i$ while authenticating (a) its type and site,
      (b) $G_i(v)$ in the exact source commitment context, and (c) a recovery
      token that will disclose this same $v$ at every direct reveal of the
      binding.  A backend may implement this by public checking, a proof plus
      escrow, trusted ideal functionality, or another explicit mechanism.
      Merely accepting an opaque unopenable handle is insufficient to establish
      a source commitment step.  In particular, an invalid or fresh
      unprepared candidate is rejected before $pc$ advances.
\item Under that sufficient alternative, a direct reveal has mode
      $\mathsf{recover}$: the owner may open first,
      but on timeout the recovery capability discloses the certified value.
      It never substitutes a different value.  Failure to provide the recovery
      capability is an unsupported deployment, not automatically a source
      quit, because source reveal is deterministic.
\item Initial-input availability loads a typed private setup supplied outside
      public code.  An initial deterministic reveal is automatic once enabled;
      setup failure is an admission failure, not a player deviation.
\item A chance capability samples exactly once from $K_i$ evaluated on the
      effective public prefix, caches the result, and publishes it atomically.
      An unavailable sampler stutters rather than substitutes a default.
\end{enumerate}

Private guards need not expose private reads.  What is required is evidence
that the accepted value satisfies the source guard in the owner's source view.
Public-guard checking is one implementation of this obligation, not the
definition of the language boundary.

The authenticated commitment context must itself be rooted in the ordered
source prefix: each public read comes from an earlier realized public site, and
each private read comes from initial setup or an earlier certified binding of
the same owner.  A certificate over an arbitrary raw snapshot is not enough.
This recursive provenance is established incrementally by the program-counter
induction; it is not an additional oracle claiming final settlement legality.

No general field of the form ``fallback is source-correct'' should be added to
each site.  Such a field would simply assume the difficult refinement theorem.
Resolution is instead justified by recognizable source syntax, such as the
optional conditional-opening pair, or by a concrete recovery or settlement
mechanism whose correspondence is proved separately.

This is one concrete mode discipline for legal bindings and deterministic
disclosure, not a claim that recovery or escrow is indispensable or that it has
been selected as the full-Core target.  It does not by itself resolve refusal
to submit a commitment: that case still requires an explicit, source-related
resolution or an admission contract.  Recoverability after acceptance is not
a complete progress or compiler theorem.
Under it, ordinary source commitments use certified recoverable bindings and
direct reveals use $\mathsf{recover}$.  Strategic disclosure failure is
represented by an explicit source decision, most simply the adjacent
optional-copy commitment and reveal: the owner chooses $\None$ or $\Some(x)$,
and the ordered machine does not pass that reveal before the choice is
materialized.  A deferred-validation source semantics could justify a
different mode, but only after its own legal executions and guard context have
been specified and related to the runtime.  A target may also offer a plain
opaque-binding mode, but the positive argument below is conditional on a
correspondence theorem for whichever mode is chosen.

\section{Why naive per-cell defaults are not sound}

Consider one owner and two source commitments
\[
 x,y : \mathsf{Option}(\mathsf{Bool}).
\]
The guard for $x$ may be the ordinary nullable guard.  The guard for $y$ allows
$y=\None$ in every view and allows $y=\Some(b)$ only if $x=\Some(b)$.  Order the
source operations as
\[
 \mathsf{commit}\ x;\quad \mathsf{commit}\ y;\quad
 \mathsf{reveal}\ y;\quad \mathsf{reveal}\ x.
\]
The source execution that binds $x=\Some(b)$ and $y=\Some(b)$ is legal.  A
target can first publish $y=\Some(b)$ and later fail to open $x$.  If timeout
rewrites the effective and public value of the original $x$ to $\None$, the
final public pair is
\[
 (x,y)=(\None,\Some(b)).
\]
No legal source execution has this result.  Rechecking only future openings
does not help, because the dependent value $y$ has already been published.
Likewise, validating $y$ against a captured original snapshot proves only that
it was legal when bound; it does not make the rewritten terminal store a legal
source store.

This failure persists even under a strong source quitting inequality.  Define
utility to be $2$ on the impossible pair $(\None,\Some(b))$, $1$ on a valid
both-present pair, and $0$ otherwise.  Every legal source quit outcome has
utility $0$, and every legal supported continuation has nonnegative utility.
The pointwise source quit condition therefore holds, while the naive target
timeout can create utility $2$.  The missing premise is operational
correspondence, not a stronger source incentive inequality.

This is only a counterexample to direct null replacement of an original source
cell while retaining already published dependent results.  It is neither a
universal compiler impossibility nor proof that any particular stronger
mechanism is necessary.  Candidate sound implementation choices include:
\begin{enumerate}
\item recover or escrow deterministic openings, so the later timeout cannot
      rewrite $x$;
\item expose a distinct source-declared optional opening $o\in
      \{\None,\Some(x)\}$ and resolve $o$ to $\None$ while retaining $x$;
\item use an explicit settlement-closure mechanism that proves all already
      public dependent effects admit one legal source realization.
\end{enumerate}
For unpublished descendants, a topological cascade of source-legal defaults
may be provable.  It cannot retract an already published descendant and hence
is not by itself a general solution.

\section{The correspondence certificate}

The target model should be expressive enough to run the counterexample above,
but the compiler theorem must refuse to claim source correspondence for its
naive settlement.  The smallest load-bearing certificate is therefore not a
per-node fallback oracle.  It is a proof, derived for a concrete lowering and
execution, of the following alternative.

\paragraph{Proposed first-resolution realization.}
For every target trace, either it realizes the ordinary source prefix, or at
the first source-relevant resolution there are
\begin{itemize}
\item a designated source decision $d$ owned by the focal player,
\item a legal terminal source execution $q$ that takes the declared resolution
      at $d$,
\item a continued legal source execution $c$ under a unilateral source policy,
\end{itemize}
such that $q$ reproduces the target's payoff-relevant public settlement, $q$
and $c$ agree on the required pre-decision source information, and all earlier
target blocks correspond to their source operations.

This certificate may be proved for an adjacent conditional opening, for an
escrow-backed deterministic reveal, or for another explicit mechanism.  It is
false for the naive reversed-reveal trace above.  Keeping it as a theorem about
the lowering avoids weakening the source language and avoids installing a
circular global-refinement assumption inside the runtime.

\section{Proposed local induction lemmas}

The following are proof targets, not currently asserted Lean facts.

\begin{enumerate}
\item \textbf{Typed block law.}  With no resolution, executing site $i$ and
      advancing $pc$ realizes exactly the corresponding heterogeneous source
      operation.  Malformed, mistyped, replayed, and off-order traffic does not
      change the realized source prefix.
\item \textbf{Binding law.}  Binding evidence reconstructs the source guard
      proof in the exact owner view.  It preserves the raw value and the
      owner's private-action recall without exposing private inputs.
\item \textbf{Opening law.}  A successful deterministic opening publishes the
      bound source value.  A conditional-opening resolution publishes the
      optional copy's legal $\None$, not a mutation of its source secret.
\item \textbf{Chance law.}  Conditional on an equal effective public prefix,
      the cached target sample has kernel $K_i$ and is sampled only once.
\item \textbf{Observation law.}  Equal source decision views induce equal
      retained target observations, including the focal player's own earlier
      choices.  Extra packet history may be retained unless a separate fixed-
      environment factorization theorem proves it irrelevant.
\item \textbf{First-resolution decomposition.}  Every stopped target trace
      either has no source-relevant resolution or has an earliest such
      resolution; the preceding ordered prefix satisfies the block laws.
\item \textbf{Resolution realization.}  A recognized resolution mechanism
      constructs the legal $q$ described above and proves public-settlement
      agreement.  This is exactly where the reversed-reveal counterexample
      blocks a naive backend.
\end{enumerate}

Together these lemmas support induction over the public program counter.  They
should be ordinary theorems consumed by the compiler proof, not fields saying
that all global executions or deviations are already correct.

\section{Strategies, environment, and causal deviations}

Fix a focal player $p$.  A strategy is an observation-local stochastic kernel
with perfect recall of $p$'s private setup and prior actions.  Fix all opponent
strategies and an adaptive environment strategy before choosing a deviation of
$p$.  The environment may react to its observations, deliver or delay
messages, and choose when to advance time.  It must not be reselected after
seeing which focal policy is under comparison.

The relevant fairness condition is timeout-relative and uniform over focal
deviations.  For a nonfocal honest site, once its prescribed honest action is
enabled and available, the fixed environment must service that action before
the site's deadline; it may not let timeout win the race.  A focal player may
withhold an action, after which the declared resolution mechanism may run at
the deadline.  Chance and automatic initial disclosure receive their declared
service.  Fairness does not require a fixed absolute schedule, and it does not
erase the environment's lawful adaptive behavior.

\paragraph{Proposed causal extraction.}
For every observation-local target deviation $\tau_p$, with the opponents and
environment fixed as above, construct a source behavioral policy $\alpha_p$
by induction over retained source decisions.  The action law at a source view
is obtained by conditioning the target law on the corresponding retained
history.  Random response tapes may be sampled ex ante.  The construction must
not choose a source action post hoc using hidden future transport, chance, or
opponent outcomes.

The induction should produce a joint law of target result $R$, continued source
execution $C$, and, on a resolution branch, legal quitting execution $Q$ and
site $d$.  Crucially, the $C$ marginal must be the complete source law induced
by the extracted policy, not merely a distribution supported by that law:
\begin{align*}
 \Law(C) &= \Law(P,\alpha_p,\sigma_{-p}),\\
 \pub(R) &= \pub(C) \quad\text{on ordinary branches},\\
 \pub(R) &= \pub(Q) \quad\text{on resolution branches},\\
 \Obs_d(Q) &= \Obs_d(C) \quad\text{on resolution branches}.
\end{align*}
An alternative formulation may use a mixture of source policies, but it must
state its exact mixed marginal and obtain the corresponding source incentive
bound.  Mere pointwise support membership is insufficient for the expected
utility inequality below.
Here $\Obs_d$ is source-defined.  The existing public-prefix predicate is a
sufficient choice when the realization proof supplies its equality.  A weaker
information-indexed incentive theorem may instead use the focal source view,
but its probability law and conditioning must also be proved by this coupling;
it cannot be assumed as a native-conditioned premise.

\section{Sufficient incentive conditions}

Let $U_p$ be a utility of the public source outcome.  The presently usable
pointwise condition is
\[
 \Obs_d(q)=\Obs_d(c)
 \quad\Longrightarrow\quad
 U_p(q)\leq U_p(c)
\]
for every legal declared-resolution execution $q$ and every supported
unilateral continuation $c$ at the same owned decision.  The existing
public-prefix quitting condition is stronger than necessary whenever the
coupling preserves a richer focal view, but it is already adequate once the
legal $q$ exists.

A genuinely weaker future condition could compare conditional expectations on
a source-defined observation cylinder $Z$:
\[
 \mathbb{E}[U_p(Q)\mid Z=z]
 \leq
 \mathbb{E}[U_p(C)\mid Z=z].
\]
It is useful only if the compiler proves the joint cylinder probabilities and
the stopping coupling.  Quantifying over all opponent restrictions and all
deterministic focal alternatives may isolate every positive-probability source
outcome and collapse this condition back to pointwise dominance.  Therefore an
admissible cylinder family must come from the retained source observation, not
from arbitrary post hoc restrictions.

Assume the ordinary branch and first-resolution coupling above, the pointwise
condition, and an honest exact outcome law.  Then, for every target deviation,
\[
 \mathbb{E}U_p(R_{\tau_p})
 \leq \mathbb{E}U_p(C_{\alpha_p})
 \leq \mathbb{E}U_p(C_{\sigma_p})+\varepsilon
 = \mathbb{E}U_p(R_{\sigma_p})+\varepsilon,
\]
where the middle inequality is source $\varepsilon$-Nash.  The first inequality
uses legal resolution realization plus quitting dominance; it is precisely the
step unavailable for the naive per-cell-default counterexample.

\section{Adequacy checklist for a full-Core lowering}

\begin{center}
\begin{tabular}{>{\raggedright\arraybackslash}p{0.23\linewidth}
                >{\raggedright\arraybackslash}p{0.69\linewidth}}
Source feature & Required target fact \\
\hline
Heterogeneous values & Typed tags and dependent stores; no homogeneous widening.\\
Private guarded choice & Exact-view binding evidence, without public disclosure of private reads.\\
Public guarded choice & Direct public validator is an admissible specialization.\\
Deterministic reveal & Under current semantics, an availability or recovery correspondence; alternatively, an explicitly redesigned deferred-validation source step.\\
Conditional opening & Preserve the original secret and resolve the fresh optional copy.\\
Public dependent chance & Evaluate on the effective public prefix, sample once, and cache.\\
Initial private setup & Separate typed input plus automatic availability; not a strategic publication.\\
Timeout & Attribute it to an explicit source-resolvable operation and prove legal settlement realization.\\
Arbitrary deviation & Observation-local causal policy extraction against opponents and environment fixed in advance.\\
Adaptive timing & Uniform timeout-relative fairness, not a policy-dependent schedule.\\
\end{tabular}
\end{center}

The model may express deployments lacking one of these capabilities.  What it
must not do is attach a source-correctness theorem to such a deployment.  This
separation keeps the target language expressive while making the compiler's
positive theorem both meaningful and falsifiable.

\end{document}
